% Book 3, Year 1, chapter 17 -- binary search trees and balanced trees.
% Provenance (never printed): sources/book3/BRIEFS.md, chapter 17.
% ledger: b1:avl-origin, b1:rb-origin, b1:llrb, b2:map-complexity
\chapter{Binary Search Trees and Balanced Trees}\label{ch:b1:balanced-trees}

The town library keeps its catalogue in a search tree, so that a title is found by a few comparisons and the titles can
be listed in alphabetical order. One year the librarian rebuilds the tree by inserting the titles in alphabetical
order, from a printed list. The new tree answers correctly, but every search now takes forty comparisons on average
instead of a handful: the tree has become a \omterm{def:b1:linked-structures:list}{linked list}. This chapter builds search trees, measures how their shape depends
on the order of insertion, and makes two kinds of trees that keep their height logarithmic whatever the order: \omterm{def:b1:balanced-trees:avl}{AVL trees}
and \omterm{def:b1:balanced-trees:rb}{red-black trees}.

\begin{recall}
A \omterm{def:b1:induction-and-termination:binary-tree}{binary tree} is empty or a root with a left and a right \omterm{def:b1:induction-and-termination:rooted-tree}{subtree}; its height is the number of edges on a longest path
from the root down, $-1$ for the empty tree (\cref{def:b1:induction-and-termination:height}). Properties of trees are
proved by \omterm{def:b1:induction-and-termination:structural}{structural induction} (\cref{def:b1:induction-and-termination:structural}). A \omterm{def:b1:hash-tables:dictionary}{dictionary} associates values with
keys (\cref{def:b1:hash-tables:dictionary}).
\end{recall}

\section{Binary search trees}\label{sec:b1:balanced-trees:bst}

\begin{definition}[Binary search tree]\label{def:b1:balanced-trees:bst}
A \emph{binary search tree}\index{binary search tree} is a \omterm{def:b1:induction-and-termination:binary-tree}{binary tree} whose nodes hold keys from an ordered set, with
the \emph{search-tree property}\index{search-tree property}: at every node, the keys of the left \omterm{def:b1:induction-and-termination:rooted-tree}{subtree} are smaller than
the node's key, and the keys of the right \omterm{def:b1:induction-and-termination:rooted-tree}{subtree} are larger.
\end{definition}

A search compares the key with the root's and goes left or right; it follows one path, so it costs $O(h)$ for a tree of
height $h$. Insertion searches, and puts the new node where the search fell off the tree.

\omcode{code/bachelor-1/17-balanced-trees/listings/bst.c}{38}{47}{The floor: the largest key at most
\texttt{key}}\label{lst:b1:balanced-trees:floor}

\begin{definition}[Traversals]\label{def:b1:balanced-trees:traversal}
The \emph{in-order traversal}\index{in-order traversal} of a \omterm{def:b1:induction-and-termination:binary-tree}{binary tree} visits the left \omterm{def:b1:induction-and-termination:rooted-tree}{subtree}, then the root, then the
right \omterm{def:b1:induction-and-termination:rooted-tree}{subtree}; the \emph{preorder}\index{preorder} visits the root first, the \emph{postorder}\index{postorder} last.
\end{definition}

\begin{proposition}[In-order is sorted]\label{prop:b1:balanced-trees:inorder}
The \omterm{def:b1:balanced-trees:traversal}{in-order traversal} of a \omterm{def:b1:balanced-trees:bst}{binary search tree} visits its keys in increasing order.
\end{proposition}

\begin{proof}
By \omterm{def:b1:induction-and-termination:structural}{structural induction}. The empty tree visits nothing. For a root $r$ with \omterm{def:b1:induction-and-termination:rooted-tree}{subtrees} $L$ and $R$, the traversal visits
the keys of $L$ in increasing order (hypothesis), all smaller than $r$ (\omterm{def:b1:balanced-trees:bst}{search-tree property}), then $r$, then the keys of
$R$ in increasing order, all larger than $r$.
\end{proof}

\begin{definition}[Ordered dictionary, successor]\label{def:b1:balanced-trees:ordered}
An \emph{ordered dictionary}\index{ordered dictionary} is a \omterm{def:b1:hash-tables:dictionary}{dictionary} whose keys are ordered, with the further
operations minimum, maximum, floor (the largest key at most $k$), ceiling, \emph{successor}\index{successor (in a search
tree)} (the smallest key larger than $k$), and the listing of the keys between two bounds.
\end{definition}

A \omterm{def:b1:hash-tables:hash-table}{hash table} cannot answer these questions without examining every key; a search tree answers each by one path. The
floor goes right whenever the node's key is small enough, remembering it as a candidate, since a larger candidate can
only be on its right.

\omoutput{code/bachelor-1/17-balanced-trees/outputs/bst_demo.txt}{Nine names, drawn sideways (root on the left, right
\omterm{def:b1:induction-and-termination:rooted-tree}{subtree} above); the deletion of a node with two children.}\label{out:b1:balanced-trees:demo}

Deletion has three cases. A \omterm{def:b1:induction-and-termination:rooted-tree}{leaf} is removed. A node with one child is replaced by that child. A node with two children
takes the key of its \omterm{def:b1:balanced-trees:ordered}{successor}, the minimum of its right \omterm{def:b1:induction-and-termination:rooted-tree}{subtree}, which has no left child, and that \omterm{def:b1:balanced-trees:ordered}{successor} is
deleted from the right \omterm{def:b1:induction-and-termination:rooted-tree}{subtree} by one of the first two cases.

\omcode{code/bachelor-1/17-balanced-trees/listings/bst.c}{57}{72}{Deletion: the two-child case moves the \omterm{def:b1:balanced-trees:ordered}{successor}
up}\label{lst:b1:balanced-trees:delete}

\begin{omfigure}
\begin{tikzpicture}[level distance=9mm, every node/.style={font=\small, draw, rounded corners=2pt, minimum
  width=9mm, inner sep=2pt}, level 1/.style={sibling distance=26mm}, level 2/.style={sibling distance=14mm},
  level 3/.style={sibling distance=10mm}, on/.style={fill=omCode!20, thick}]
\node[on] {Mia}
  child {node[on] {Eli} child {node {Bea}} child {node[on] {Ivy} child {node {Gus}} child {node[on] {Kai}}}}
  child {node {Tom} child {node {Pia}} child {node {Zoe}}};
\begin{scope}[xshift=7.2cm]
\node {Mia}
  child {node[fill=omThm!20] {Gus} child {node {Bea}} child {node {Ivy} child[missing] child {node {Kai}}}}
  child {node {Tom} child {node {Pia}} child {node {Zoe}}};
\end{scope}
\end{tikzpicture}
\par\smallskip{\small Left: the tree of \cref{out:b1:balanced-trees:demo} and the path of the search for the floor of
\code{Jo} (shaded): Mia, Eli, Ivy and Kai; the answer is Ivy, the last node where the search went right. Right: after
deleting Eli, its \omterm{def:b1:balanced-trees:ordered}{successor} Gus has taken its place.}
\end{omfigure}

\section{Height bounds}\label{sec:b1:balanced-trees:height}

Every operation of the previous section costs $\Theta(h)$ in the worst case. The height depends on the order of the
insertions, not only on the keys.

\begin{proposition}[Height of a binary tree]\label{prop:b1:balanced-trees:height}
A \omterm{def:b1:induction-and-termination:binary-tree}{binary tree} with $n \geq 1$ nodes has height $h$ with $\lfloor \log_2 n \rfloor \leq h \leq n - 1$.
\end{proposition}

\begin{proof}
A longest path holds $h + 1$ nodes, all distinct, so $n \geq h + 1$. Depth $d$ holds at most $2^d$ nodes, so $n \leq 1 +
2 + \dots + 2^h = 2^{h+1} - 1 < 2^{h+1}$, and $\log_2 n < h + 1$, that is $\lfloor \log_2 n \rfloor \leq h$.
\end{proof}

\begin{definition}[Degenerate tree]\label{def:b1:balanced-trees:degenerate}
A tree is \emph{degenerate}\index{degenerate tree} when every node has at most one child: its height is $n - 1$, and it
is a \omterm{def:b1:linked-structures:list}{linked list}.
\end{definition}

Inserting keys in increasing order builds a \omterm{def:b1:balanced-trees:degenerate}{degenerate tree}: each new key is larger than all the others and goes to the
right of the rightmost node. The town's 80 titles show the three regimes.

\omoutput{code/bachelor-1/17-balanced-trees/outputs/heights.txt}{Heights of trees built from the first $n$ titles.
The last three lines are used by the problem.}\label{out:b1:balanced-trees:heights}

\begin{omfigure}
\begin{tikzpicture}
\begin{axis}[width=11cm, height=6cm, xmin=0, xmax=82, ymin=0, ymax=26, xlabel={number of titles $n$},
  ylabel={height}, legend pos=north east, legend style={font=\small}, tick label style={font=\small}]
\addplot[thick, omThm, dashed] table[x=n, y=sorted] {code/bachelor-1/17-balanced-trees/figures/heights.dat};
\addlegendentry{plain tree, alphabetical order}
\addplot[thick, omThm, mark=*, mark size=1pt] table[x=n, y=shuffled]
  {code/bachelor-1/17-balanced-trees/figures/heights.dat}; \addlegendentry{plain tree, random order (mean)}
\addplot[thick, omDef, mark=square*, mark size=1pt] table[x=n, y=avl]
  {code/bachelor-1/17-balanced-trees/figures/heights.dat}; \addlegendentry{AVL tree}
\addplot[thick, black, densely dotted] table[x=n, y=lg] {code/bachelor-1/17-balanced-trees/figures/heights.dat};
\addlegendentry{$\lfloor \log_2 n \rfloor$}
\end{axis}
\end{tikzpicture}
\par\smallskip{\small Heights from \cref{out:b1:balanced-trees:heights}. The alphabetical order climbs as $n - 1$ and
leaves the chart at $n = 27$; random orders grow slowly; the \omterm{def:b1:balanced-trees:avl}{AVL tree} stays within one or two of the lower bound.}
\end{omfigure}

A random order of insertion gives, on average, a height of order $\log n$; the average depth of a node is about $2 \ln n
\approx 1.39 \log_2 n$, a result of the second year. Input does not come in random order: sorted files, increasing
identifiers and timestamps are common. Balanced trees guarantee $O(\log n)$ for every order, by repairing the tree after
each insertion with \omterm{def:b1:balanced-trees:rotation}{rotations}.

\section{Rotations}\label{sec:b1:balanced-trees:rotations}

\begin{definition}[Rotation]\label{def:b1:balanced-trees:rotation}
A \emph{right rotation}\index{rotation} at a node $y$ with left child $x$ makes $x$ the root of the \omterm{def:b1:induction-and-termination:rooted-tree}{subtree}, $y$ its right
child, and the former right \omterm{def:b1:induction-and-termination:rooted-tree}{subtree} of $x$ the left \omterm{def:b1:induction-and-termination:rooted-tree}{subtree} of $y$. A \emph{left rotation} is the inverse operation.
\end{definition}

\begin{omfigure}
\begin{tikzpicture}[level distance=10mm, n/.style={draw, circle, minimum size=7mm, font=\small},
  t/.style={draw, isosceles triangle, shape border rotate=90, minimum height=8mm, minimum width=8mm, inner sep=0pt,
  font=\small, fill=omCode!10}, level 1/.style={sibling distance=20mm}, level 2/.style={sibling distance=11mm},
  >={Stealth[length=2mm]}]
\node[n] {$y$} child {node[n] {$x$} child {node[t] {$A$}} child {node[t, fill=omThm!20] {$B$}}}
  child {node[t] {$C$}};
\draw[->, thick] (2.2,-1.2) -- node[above, font=\small] {right at $y$} (4.0,-1.2);
\draw[<-, thick] (2.2,-1.6) -- node[below, font=\small] {left at $x$} (4.0,-1.6);
\begin{scope}[xshift=6.2cm]
\node[n] {$x$} child {node[t] {$A$}} child {node[n] {$y$} child {node[t, fill=omThm!20] {$B$}}
  child {node[t] {$C$}}};
\end{scope}
\end{tikzpicture}
\par\smallskip{\small A \omterm{def:b1:balanced-trees:rotation}{rotation} changes three links. The in-order sequence is $A, x, B, y, C$ on both sides; the \omterm{def:b1:induction-and-termination:rooted-tree}{subtree}
$B$ (shaded) changes parent; $A$ moves up a level and $C$ down a level.}
\end{omfigure}

\begin{proposition}[Rotations preserve the order]\label{prop:b1:balanced-trees:rotation}
A \omterm{def:b1:balanced-trees:rotation}{rotation} of a \omterm{def:b1:balanced-trees:bst}{binary search tree} yields a \omterm{def:b1:balanced-trees:bst}{binary search tree} with the same keys and the same in-order sequence.
\end{proposition}

\begin{proof}
Before a \omterm{def:b1:balanced-trees:rotation}{right rotation}, the in-order sequence of the \omterm{def:b1:induction-and-termination:rooted-tree}{subtree} is $A, x, B, y, C$, where $A$, $B$ and $C$ stand for the
sequences of the three \omterm{def:b1:induction-and-termination:rooted-tree}{subtrees}: $x$ is the left child of $y$, $A$ and $B$ its \omterm{def:b1:induction-and-termination:rooted-tree}{subtrees}. After it, $x$ is the root with
left \omterm{def:b1:induction-and-termination:rooted-tree}{subtree} $A$ and right \omterm{def:b1:induction-and-termination:rooted-tree}{subtree} $y$, whose \omterm{def:b1:induction-and-termination:rooted-tree}{subtrees} are $B$ and $C$: the sequence is again $A, x, B, y, C$. The rest of
the tree, outside the \omterm{def:b1:induction-and-termination:rooted-tree}{subtree}, is unchanged, and the \omterm{def:b1:induction-and-termination:rooted-tree}{subtree}'s set of keys too, so the property holds everywhere.
\end{proof}

\omcode{code/bachelor-1/17-balanced-trees/listings/avl.c}{18}{25}{A \omterm{def:b1:balanced-trees:rotation}{right rotation}, with the stored heights
updated}\label{lst:b1:balanced-trees:rotate}

\section{AVL trees}\label{sec:b1:balanced-trees:avl}

\begin{definition}[Balance factor, AVL tree]\label{def:b1:balanced-trees:avl}
The \emph{balance factor}\index{balance factor} of a node is the height of its right \omterm{def:b1:induction-and-termination:rooted-tree}{subtree} minus the height of its left
\omterm{def:b1:induction-and-termination:rooted-tree}{subtree}. An \emph{AVL tree}\index{AVL tree} is a \omterm{def:b1:balanced-trees:bst}{binary search tree} in which every node has balance factor $-1$, 0 or 1.
\end{definition}

\omterm{def:b1:balanced-trees:avl}{AVL trees} take their name from G. M. Adelson-Velsky and E. M. Landis, who published them in 1962. Each node stores its
height. After an insertion, the heights change only on the path from the new node up to the root, and a node on it may
reach \omterm{def:b1:balanced-trees:avl}{balance factor} $\pm 2$. Four cases repair it, named after the path from the unbalanced node to the new key: a
single \omterm{def:b1:balanced-trees:rotation}{rotation} when the path goes twice the same way (LL or RR), a double \omterm{def:b1:balanced-trees:rotation}{rotation} when it turns (LR or RL), which
first rotates the child the other way.

\omcode{code/bachelor-1/17-balanced-trees/listings/avl.c}{36}{49}{Rebalancing on the way back up}\label{lst:b1:balanced-trees:rebalance}

\begin{omfigure}
\begin{tikzpicture}[level distance=9mm, n/.style={draw, circle, minimum size=6mm, font=\small, inner sep=1pt},
  new/.style={n, fill=omThm!25}, level 1/.style={sibling distance=10mm}, >={Stealth[length=2mm]}]
\node[font=\small\bfseries] at (0.6,0.8) {LL: single rotation};
\node[n] (a) at (0,0) {3} child {node[n] {2} child {node[new] {1}} child[missing]} child[missing];
\draw[->, thick] (1.0,-0.9) -- (1.9,-0.9);
\node[n] at (3.0,0) {2} child {node[n] {1}} child {node[n] {3}};
\begin{scope}[xshift=7cm]
\node[font=\small\bfseries] at (1.9,0.8) {LR: double rotation};
\node[n] at (0,0) {3} child {node[n] {1} child[missing] child {node[new] {2}}} child[missing];
\draw[->, thick] (0.9,-0.9) -- node[above, font=\scriptsize] {left at 1} (2.1,-0.9);
\node[n] at (3.0,0) {3} child {node[n] {2} child {node[n] {1}} child[missing]} child[missing];
\draw[->, thick] (3.6,-0.9) -- node[above, font=\scriptsize] {right at 3} (4.8,-0.9);
\node[n] at (5.8,0) {2} child {node[n] {1}} child {node[n] {3}};
\end{scope}
\end{tikzpicture}
\par\smallskip{\small Two of the four rebalancing cases, on three keys; the new key is shaded. The node 3 has \omterm{def:b1:balanced-trees:avl}{balance
factor} $-2$. When the new key lies left of its left child (LL), one \omterm{def:b1:balanced-trees:rotation}{right rotation} at 3 restores the balance. When it
lies to the right of its left child (LR), a left \omterm{def:b1:balanced-trees:rotation}{rotation} at 1 turns the case into LL. RR and RL are the mirror images.}
\end{omfigure}

\omoutput{code/bachelor-1/17-balanced-trees/outputs/avl_demo.txt}{Increasing keys, then two keys that need double
\omterm{def:b1:balanced-trees:rotation}{rotations}; \omterm{def:b1:balanced-trees:avl}{balance factors} in parentheses.}\label{out:b1:balanced-trees:avl-demo}

Ten increasing keys, which built a \omterm{def:b1:balanced-trees:degenerate}{degenerate tree} of height 9 in a plain search tree, give an \omterm{def:b1:balanced-trees:avl}{AVL tree} of height 3. The
cost of an insertion is the descent, $O(h)$, plus constant work at each node of the way back up: at most one \omterm{def:b1:balanced-trees:rotation}{rotation}
event happens per insertion (\cref{exo:b1:balanced-trees:11}). The height is guaranteed by counting the smallest trees.

\begin{theorem}[Height of an AVL tree]\label{thm:b1:balanced-trees:avl-height}
An \omterm{def:b1:balanced-trees:avl}{AVL tree} with $n$ nodes has height $h \leq \log_\varphi (n + 1) - 1 < 1.45 \log_2 (n + 2)$, where $\varphi = (1 +
\sqrt 5)/2$.
\end{theorem}

\begin{proof}
Let $N(h)$ be the smallest number of nodes of an \omterm{def:b1:balanced-trees:avl}{AVL tree} of height $h$: $N(0) = 1$, $N(1) = 2$. A smallest tree of height
$h \geq 2$ has one \omterm{def:b1:induction-and-termination:rooted-tree}{subtree} of height $h - 1$ and, since the \omterm{def:b1:balanced-trees:avl}{balance factor} is at most 1 in absolute value and the tree
is smallest, the other of height $h - 2$, both smallest themselves: $N(h) = N(h - 1) + N(h - 2) + 1$. With the Fibonacci
numbers $F_0 = 0$, $F_1 = 1$, $F_{k+2} = F_{k+1} + F_k$, induction gives $N(h) = F_{h+3} - 1$, and $F_k \geq
\varphi^{k-2}$ for $k \geq 1$ (problem, Part IV). So $n \geq N(h) \geq \varphi^{h+1} - 1$, and $h + 1 \leq \log_\varphi
(n + 1)$. Finally $1/\log_2 \varphi \approx 1.4404$.
\end{proof}

\omoutput{code/bachelor-1/17-balanced-trees/outputs/minimal_avl.txt}{Minimal \omterm{def:b1:balanced-trees:avl}{AVL trees}, built as trees and
counted.}\label{out:b1:balanced-trees:minimal}

\section{Red-black trees}\label{sec:b1:balanced-trees:rb}

\begin{definition}[Red-black tree, black height]\label{def:b1:balanced-trees:rb}
A \emph{red-black tree}\index{red-black tree} is a \omterm{def:b1:balanced-trees:bst}{binary search tree} whose nodes are coloured red or black, such that
the root is black, a red node has no red child, and every path from a node down to an empty \omterm{def:b1:induction-and-termination:rooted-tree}{subtree} contains the same
number of black nodes, the \emph{black height}\index{black height} of the node. A \emph{left-leaning red-black
tree}\index{left-leaning red-black tree} is a red-black tree in which every red node is a left child.
\end{definition}

\omterm{def:b1:balanced-trees:rb}{Red-black trees} were introduced by L. Guibas and R. Sedgewick in 1978; the left-leaning variant is
Sedgewick's. A red node and its parent can be read as one node with two keys: a \omterm{def:b1:balanced-trees:rb}{left-leaning red-black tree} is then a
\emph{2-3 tree}, whose nodes hold one or two keys and whose empty \omterm{def:b1:induction-and-termination:rooted-tree}{subtrees} are all at the same depth. Three local repairs,
on the way back up after inserting a red node, keep the colour rules: a red right link is rotated to the left, two red
links in a row are rotated right, and a node with two red children passes its red up by a colour flip.

\omcode{code/bachelor-1/17-balanced-trees/listings/llrb.c}{32}{47}{Insertion in a \omterm{def:b1:balanced-trees:rb}{left-leaning red-black
tree}}\label{lst:b1:balanced-trees:llrb}

\omoutput{code/bachelor-1/17-balanced-trees/outputs/llrb_demo.txt}{The same ten increasing keys; red nodes are
starred.}\label{out:b1:balanced-trees:llrb-demo}

\begin{omfigure}
\begin{tikzpicture}[level distance=9mm, n/.style={draw, circle, minimum size=6mm, font=\small, inner sep=1pt,
  fill=black!75, text=white}, r/.style={draw=omThm, circle, minimum size=6mm, font=\small, inner sep=1pt,
  fill=omThm!25}, level 1/.style={sibling distance=36mm}, level 2/.style={sibling distance=20mm},
  level 3/.style={sibling distance=9mm}]
\node[n] {D}
  child {node[n] {B} child {node[n] {A}} child {node[n] {C}}}
  child {node[n] {H} child {node[r] {F} child {node[n] {E}} child {node[n] {G}}}
    child {node[n] {J} child {node[r] {I}} child[missing]}};
\begin{scope}[xshift=8.3cm, b/.style={draw, rounded corners=2pt, font=\small, minimum height=6mm, inner sep=3pt},
  level 1/.style={sibling distance=30mm}, level 2/.style={sibling distance=12mm}]
\node[b] {D}
  child {node[b] {B} child {node[b] {A}} child {node[b] {C}}}
  child {node[b] {F\enspace H} child {node[b] {E}} child {node[b] {G}} child {node[b] {I\enspace J}}};
\end{scope}
\end{tikzpicture}
\par\smallskip{\small Left: the \omterm{def:b1:balanced-trees:rb}{left-leaning red-black tree} of \cref{out:b1:balanced-trees:llrb-demo}, red nodes shaded;
every path from the root to an empty \omterm{def:b1:induction-and-termination:rooted-tree}{subtree} meets three black nodes. Right: the same tree read as a 2-3 tree, each red
node merged with its parent: all leaves at depth 2.}
\end{omfigure}

\begin{theorem}[Height of a red-black tree]\label{thm:b1:balanced-trees:rb-height}
A \omterm{def:b1:balanced-trees:rb}{red-black tree} with $n$ nodes has height $h$ with $h + 1 \leq 2 \log_2 (n + 1)$.
\end{theorem}

The proof (\cref{exo:b1:balanced-trees:10}) counts at least $2^b - 1$ nodes below a node of \omterm{def:b1:balanced-trees:rb}{black height} $b$, and at least
one black node out of two on any path. The bound is weaker than the \omterm{def:b1:balanced-trees:avl}{AVL tree}'s, $2$ against $1.44$, but the repairs are
simpler, and deletion needs at most a constant number of \omterm{def:b1:balanced-trees:rotation}{rotations} in the classic version. The standard libraries of
C++ and Java implement their ordered maps as \omterm{def:b1:balanced-trees:rb}{red-black trees}; the left-leaning variant shown here is shorter to write,
and its deletion, told but not written, is shorter too.

\section{Serialising a tree}\label{sec:b1:balanced-trees:serialisation}

\begin{definition}[Serialisation]\label{def:b1:balanced-trees:serialisation}
The \emph{serialisation}\index{serialisation (of a data structure)} of a data structure is a sequence of bytes, or of
characters, from which the structure can be rebuilt, in another run or on another machine; \omterm{def:b1:pointers-and-arrays:pointer}{pointers} are not part of it,
since they are only meaningful in one address space.
\end{definition}

A \omterm{def:b1:balanced-trees:traversal}{preorder} traversal that writes a marker for each empty \omterm{def:b1:induction-and-termination:rooted-tree}{subtree} describes the shape completely: reading it back, the
first item is the root, then comes the left \omterm{def:b1:induction-and-termination:rooted-tree}{subtree}, written in the same way, then the right one.

\omcode{code/bachelor-1/17-balanced-trees/listings/serialise.c}{9}{31}{Writing and reading a tree in \omterm{def:b1:balanced-trees:traversal}{preorder} with
markers}\label{lst:b1:balanced-trees:serialise}

\omoutput{code/bachelor-1/17-balanced-trees/outputs/serialise.txt}{A tree written to a file and read
back.}\label{out:b1:balanced-trees:serialise}

\begin{inpractice}
For a search tree, the sorted list of keys is enough to rebuild a balanced tree in linear time
(\cref{exo:b1:balanced-trees:12}), and it is the better file format: it does not depend on the shape, which is an
accident of the insertion order. The \omterm{def:b1:balanced-trees:traversal}{preorder} with markers is needed when the shape itself is the information, as in an
expression tree or a decision tree.
\end{inpractice}

\section{Exercises}\label{sec:b1:balanced-trees:exercises}

\begin{exercise}[$\star$]\label{exo:b1:balanced-trees:1}
The keys 50, 30, 70, 20, 40, 60, 80, 35 are inserted in this order into an empty search tree. Give the in-order, the
\omterm{def:b1:balanced-trees:traversal}{preorder} and the \omterm{def:b1:balanced-trees:traversal}{postorder} sequences of the tree.
\end{exercise}

\begin{exercise}[$\star$]\label{exo:b1:balanced-trees:2}
Draw the search tree obtained by inserting 8, 3, 10, 1, 6, 14, 4, 7, 13, and give its height. Then delete 3 and draw the
tree again.
\end{exercise}

\begin{exercise}[$\star$]\label{exo:b1:balanced-trees:3}
Insert 5, 3, 8, 2, 4, 1 into a plain search tree and give the \omterm{def:b1:balanced-trees:avl}{balance factor} of every node. Is it an \omterm{def:b1:balanced-trees:avl}{AVL tree}? Which
\omterm{def:b1:balanced-trees:rotation}{rotation} would make it one?
\end{exercise}

\begin{exercise}[$\star$]\label{exo:b1:balanced-trees:4}
Insert S, E, A, R into an empty \omterm{def:b1:balanced-trees:rb}{left-leaning red-black tree}, and give the tree and the colours after each insertion.
\end{exercise}

\begin{exercise}[$\star\star$]\label{exo:b1:balanced-trees:5}
This deletion lets one child take the place of the deleted node. Explain the output and the memory it leaks.
\omcode{code/bachelor-1/17-balanced-trees/listings/exo_delete_bug.c}{8}{19}{A shortcut}
\omoutput{code/bachelor-1/17-balanced-trees/outputs/exo_delete_bug.txt}{Deleting the root of seven names.}
\end{exercise}

\begin{exercise}[$\star\star$]\label{exo:b1:balanced-trees:6}
The \omterm{def:b1:balanced-trees:rotation}{rotation} is correct, but its caller is not. The tree held Mia, Eli, Tom, Bea, Ivy and Abe. Explain the output and
repair the call.
\omcode{code/bachelor-1/17-balanced-trees/listings/exo_rotate_bug.c}{6}{11}{A \omterm{def:b1:balanced-trees:rotation}{right rotation}}
\omoutput[2-2]{code/bachelor-1/17-balanced-trees/outputs/exo_rotate_bug.txt}{After \code{rotate_right(t->left);}.}
\end{exercise}

\begin{exercise}[$\star\star$]\label{exo:b1:balanced-trees:7}
Write \code{range_count(t, lo, hi)}, the number of keys $k$ with $\code{lo} \leq k \leq \code{hi}$, in time $O(h + r)$
where $r$ is the answer, and test it.
\end{exercise}

\begin{exercise}[$\star\star$]\label{exo:b1:balanced-trees:8}
A function checks that a tree is a search tree by comparing each node with its two children. Give a tree that it
accepts wrongly, then write a correct check that passes bounds down the tree.
\end{exercise}

\begin{exercise}[$\star\star$]\label{exo:b1:balanced-trees:9}
Count the shapes of search trees with $n$ nodes for $n \leq 5$, first by a program that inserts every permutation of
$n$ keys and collects the shapes, then by the recurrence $C_n = \sum_{i=0}^{n-1} C_i C_{n-1-i}$. Justify the recurrence.
\end{exercise}

\begin{exercise}[$\star\star\star$]\label{exo:b1:balanced-trees:10}
Prove \cref{thm:b1:balanced-trees:rb-height}: show that a \omterm{def:b1:induction-and-termination:rooted-tree}{subtree} whose root has \omterm{def:b1:balanced-trees:rb}{black height} $b$ has at least $2^b - 1$
nodes, then that $b \geq (h + 1)/2$ at the root.
\end{exercise}

\begin{exercise}[$\star\star\star$]\label{exo:b1:balanced-trees:11}
Prove that an AVL insertion performs at most one single or double \omterm{def:b1:balanced-trees:rotation}{rotation}, by showing that the \omterm{def:b1:balanced-trees:rotation}{rotation} gives the
rebalanced \omterm{def:b1:induction-and-termination:rooted-tree}{subtree} the height it had before the insertion.
\end{exercise}

\begin{exercise}[$\star\star\star$]\label{exo:b1:balanced-trees:12}
Write a function that builds a search tree of minimal height from $n$ sorted keys in time $O(n)$, and prove that its
height is $\lfloor \log_2 n \rfloor$.
\end{exercise}

\section{Problem: Indexing the library catalogue}\label{sec:b1:balanced-trees:problem}

\begin{problem}[{Weekend problem --- Indexing the library catalogue}]\label{pb:b1:balanced-trees:1}
The town library has 80 titles, numbered in its catalogue. \Cref{out:b1:balanced-trees:heights} gives the heights of
the trees built from the first $n$ titles by inserting them in catalogue order, in alphabetical order and in 1000 random
orders, into a plain search tree, and in catalogue order into an \omterm{def:b1:balanced-trees:avl}{AVL tree} and a \omterm{def:b1:balanced-trees:rb}{left-leaning red-black tree}; its last
three lines count the \omterm{def:b1:balanced-trees:rotation}{rotations} and give the \omterm{def:b1:balanced-trees:rb}{black height}. \Cref{out:b1:balanced-trees:minimal} counts minimal \omterm{def:b1:balanced-trees:avl}{AVL
trees}.

\medskip\noindent\textbf{Part I --- Insertion order.}
\begin{enumerate}
\item Why is the height of the plain tree exactly $n - 1$ in the alphabetical column?
\item In that tree, how many comparisons does a successful search make on average, for $n = 80$?
\item How many comparisons does building it take?
\item Compare the catalogue column with the random orders. What does it suggest about the catalogue numbers?
\item Compare the heights at $n = 80$ with the lower bound $\lfloor \log_2 80 \rfloor$.
\item A library adds its new titles each year in alphabetical order. Give two ways to keep its searches fast.
\end{enumerate}

\medskip\noindent\textbf{Part II --- AVL insertion.}
\begin{enumerate}[resume]
\item The alphabetical order needs 73 single \omterm{def:b1:balanced-trees:rotation}{rotations} and no double \omterm{def:b1:balanced-trees:rotation}{rotation}. Why can increasing keys never need a
double \omterm{def:b1:balanced-trees:rotation}{rotation}?
\item In catalogue order, how many of the 80 insertions rotate? What does an insertion cost in all?
\item Insert 3, 1, 2 into an empty \omterm{def:b1:balanced-trees:avl}{AVL tree}: which case arises, and what are the two \omterm{def:b1:balanced-trees:rotation}{rotations}?
\item Why must the stored heights be updated on the way back up, and why does each update cost $O(1)$?
\item The \omterm{def:b1:balanced-trees:avl}{AVL tree} of the 80 titles in catalogue order has height 7. Is a taller \omterm{def:b1:balanced-trees:avl}{AVL tree} with 80 nodes possible?
\item A deletion can need a \omterm{def:b1:balanced-trees:rotation}{rotation} at every level of the path. Why does the argument of \cref{exo:b1:balanced-trees:11}
fail for deletions?
\end{enumerate}

\medskip\noindent\textbf{Part III --- Colours and black height.}
\begin{enumerate}[resume]
\item The \omterm{def:b1:balanced-trees:rb}{red-black tree} of the 80 titles has height 7 and \omterm{def:b1:balanced-trees:rb}{black height} 5. Check \cref{thm:b1:balanced-trees:rb-height}
and the bound $n \geq 2^b - 1$.
\item In \cref{out:b1:balanced-trees:llrb-demo}, list the red nodes and the nodes of the corresponding 2-3 tree.
\item Show that a \omterm{def:b1:balanced-trees:rb}{left-leaning red-black tree} of \omterm{def:b1:balanced-trees:rb}{black height} $b$ has between $2^b - 1$ and $3^b - 1$ keys.
\item Which \omterm{def:b1:balanced-trees:rb}{black heights} are possible for 80 keys?
\item Why does a new key enter the tree as a red node?
\item Compare the guarantees and the \omterm{def:b1:balanced-trees:rotation}{rotation} counts of AVL and \omterm{def:b1:balanced-trees:rb}{red-black trees}, and say which suits a \omterm{def:b1:hash-tables:dictionary}{dictionary} with
many deletions.
\end{enumerate}

\medskip\noindent\textbf{Part IV --- Minimal AVL trees.}
\begin{enumerate}[resume]
\item Prove $N(h) = N(h - 1) + N(h - 2) + 1$ for $h \geq 2$, with $N(0) = 1$ and $N(1) = 2$.
\item Prove $N(h) = F_{h+3} - 1$ by induction.
\item Prove $F_k \geq \varphi^{k-2}$ for $k \geq 1$, using $\varphi^2 = \varphi + 1$.
\item Deduce $h \leq \log_\varphi (n + 1) - 1$ for every \omterm{def:b1:balanced-trees:avl}{AVL tree} with $n$ nodes.
\item Compute this bound for $n = 80$, and compare with the largest height of \cref{out:b1:balanced-trees:minimal} and
with the measured AVL height.
\item State the named result of the problem, \emph{minimal \omterm{def:b1:balanced-trees:avl}{AVL trees} are Fibonacci trees}, with the height bound it
gives for the 80 titles.
\end{enumerate}
\end{problem}
